\section{Step 63 Connect load, filter, encode and save}\label{pcm-wav-s63}
Select the pipeline recipe and create the real trace input for an ordinary IDE run, or use the grader's declared in-memory file. The caller adds the parsed PCM offset to source, filters into samples, encodes into destination and saves the encoder's returned length. Stop at \texttt{results\_ready} before publication. Compare the output with the worked filtered sample codes and inspect the separate clip total.

Use \texttt{s63-1} and stop at \texttt{results\_ready}.

\paragraph{Observe} The input is eighty-two bytes, output sixty bytes, and clips zero. The filtered PCM codes are 1500, minus 1500, 1500, 500, 500, 500, 4596 and minus 3596.

\paragraph{Check} Which address does the pipeline pass to filtering?
\begin{itemize}
  \item[A.] Source base plus the accepted PCM offset.
  \item[B.] Source base plus forty-four unconditionally.
  \item[C.] The descriptor address.
\end{itemize}
